% ===== §15b: the interweaving of drive in the couple =====
\section{The Interweaving of Drive in the Couple}
\label{p26:sec:interweaving}

\noindent
This section establishes what distinguishes the intimate dyad from the general structure of the desiring subject. The aim is to show that in intimate relation the fantasy through which each subject relates to its object-cause is not merely doubled, run in parallel by two private subjects, but interwoven, each entering the place of the other's object-cause, and that this interweaving is what makes a couple's shared world both singular to them and unable to be carried out of the bond. The section states the orthodox structure of fantasy, gives the interweaving that intimate relation adds to it, draws the consequence for the non-transferability of a shared world, and states the claim.

\subsection*{The orthodox structure and what intimacy adds}

Orthodox psychoanalysis writes the subject's relation to its object in the formula of fantasy, the barred subject in relation to the objet petit a, the object-cause of desire that gives the lack a form and lets desire run \parencite{lacan1977xi}. Taken at the level of the single subject, this structure is complete in itself, and one might expect that a relation between two subjects would be the placing of two such structures side by side, each partner running a private fantasy, each turning around a private object-cause, the two merely adjacent. This is not what intimate relation is, and the difference is the seed of everything distinctive about the dyad.

In intimate relation the structure is not doubled side by side but interwoven. Each becomes implicated in the other's fantasy and enters the place of the other's object-cause, so that the object around which each one's desire turns is no longer wholly private but is bound up with, and partly occupied by, the other. The other is not merely a person one desires; the other comes to stand at the very place of the object-cause, the void around which one's drive circles, so that one's own most intimate relation to lack now passes through this particular person. Correspondingly one comes to occupy that place for them. The two voids do not merge, for a void cannot merge with a void; they are interwoven, each partner's circling now running through the other's, the two circuits linked at the place where each one's manque has taken the other as its figure.

\subsection*{The consequence for a shared world}

From this interweaving follows the character of the world a couple builds, its idioms, its private jokes, its rites, the whole texture of a shared life. This world is the sediment of the interwoven fantasy. A private joke is funny to exactly two people because its charge is anchored in a shared coordinate of desire that only those two occupy; carried to a third, it is not diminished but simply absent, because the third does not stand in the interweaving that gives it its charge. The same holds of every element of a couple's culture. Its meaning is not lodged in the words or the gestures, which can be copied, but in the interweaving of drive that the words and gestures came to figure, which cannot. This is why a shared world cannot be handed over. What could be transferred, if the outer form were copied, is a shell without its charge, the words of the joke without the interwoven desire that made them funny. The non-transferability of an intimate culture is not an accident of privacy but a consequence of where its value is lodged, in an interweaving that does not exist outside the relation and cannot be carried out of it.

\begin{claim}[The interweaving of drive]
In intimate relation the fantasy of each subject is not doubled in parallel but interwoven, each partner coming to occupy the place of the other's object-cause, so that each one's circling runs through the other's. The shared world a couple builds is the sediment of this interweaving, and its non-transferability follows: what can be copied is the outer form, but the charge is lodged in an interweaving of drive that exists nowhere outside the bond.
\end{claim}

\subsection*{The interweaving and the room}

The interweaving deepens, rather than softens, the demand of the previous section that each make room for the other's core. Because each partner has come to occupy the place of the other's object-cause, the temptation to close the other's gap is stronger here than anywhere, for the other now stands at the very center one's own desire turns around, and the wish to possess that center, to have it wholly and be assured of it, presses hardest exactly where the interweaving is deepest. The account of Part Two holds all the more firmly for that. To close the other's gap, even here, above all here, would be to extinguish the desire that made the other the living center of one's own circling, and so to collapse the interweaving into a dead thing possessed. The interweaving is kept alive by the same refusal of closure that the whole argument has traced, and the room each makes for the other's unsymbolizable core is, in the interwoven dyad, the room each makes for the living center of their own desire as well.
